\section{Experimental protocol}
\label{sec:setup}

\subsection{Cohorts}
\label{sec:cohorts}

\paragraph{Sleep-EDF}
The first cohort is the Sleep-EDF expanded database \citep{goldberger2000}: 20
healthy adults, 39 overnight polysomnograms, one EEG derivation (Fpz-Cz)
recorded at 100~Hz and scored in 30~s epochs under the R\&K rules
\citep{berry2012manual}. Wake, N1, N2, N3 and REM are retained; movement and
unscoreable epochs are dropped. Signals are band-pass filtered between 0.5 and
30~Hz, resampled to 100~Hz and standardised per recording, so that the raw mean
and variance of a night carry no information about the patient into the network.

\paragraph{HMC}
The second cohort is the Haaglanden Medisch Centrum sleep staging database
\citep{hmc2022}: 151 overnight recordings from 151 patients, one EEG channel
(C4-M1) recorded at 256~Hz and scored in 30~s epochs under the AASM rules. The
same five classes are retained, the signal is band-pass filtered between 0.5
and 30~Hz, resampled to 100~Hz and standardised per recording. The processed
cohort contains 137{,}243 labeled epochs, between 470 and 1{,}311 per
recording. It differs from the first cohort in population, montage, equipment,
raw sampling rate and scoring manual, which is what makes it a test of whether
a statement about patients survives a change of cohort.

\subsection{Backbones}

Two convolutional backbones are used. EEGNet \citep{lawhern2018}, adapted to a
single channel, has $1{,}445$ parameters; SleepCNN is a four-block VGG-style
one-dimensional network with 256 channels in its last block and
$2.24\times10^{5}$ parameters. Both were trained with batch normalization; a
group-normalized copy of each was trained separately for the compatibility
experiment. Architectures, optimization and splits are identical between the
batch and group copies except for the normalization layer.

\subsection{Arms}

All arms are evaluated on the same checkpoint, the same split and the same test
epochs; only the inference-time transformation of the penultimate
representation differs.

\begin{description}
  \item[BN.] The trained checkpoint with its cohort statistics, no correction.
  \item[Marginal.] The post-hoc representation-level correction
  $z-\wh\delta_{\cR}$, with $\wh\delta_{\cR}$ the difference between the
  recording's mean representation and the training cohort's. It is the
  representation-level counterpart of AdaBN and the operation SCSN
  generalises; unlike AdaBN it applies to a group-normalized checkpoint,
  because it reads only the penultimate features.
  \item[AdaBN.] The same checkpoint with the statistics of every
  normalization layer replaced by those of the recording being scored
  \citep{li2018adabn}. This is a nonlinear, layer-wise operation; the
  representation-level marginal correction of
  Proposition~\ref{prop:marginal} is the post-hoc arm
  \emph{Marginal} below, which removes one vector per recording and is the
  operation SCSN generalises.
  \item[SCSN.] Stage-conditional shift normalization of
  Equations~\eqref{eq:estimator}--\eqref{eq:scsn}, $\lambda=1$,
  $\lambda_s=10$. The posterior of the uncorrected model is used.
  \item[AdaBN+SCSN.] The stacked arm: layerwise AdaBN first, then SCSN in the
  adapted feature space, with prototypes estimated on source subjects that
  have each been adapted to their own recording in the same way. It tests
  whether the two corrections remove the same component or different ones.
  \item[SCSN-$\lambda$.] The interpolated correction
  $z-\delta_{\cR}-\lambda(\sum_k p_k\hat\delta_{\cR,k}-\delta_{\cR})$ for
  $\lambda\in\{0.25,0.5,0.75\}$; $\lambda=0$ is the post-hoc marginal
  correction and $\lambda=1$ is SCSN.
  \item[SCSN-$r$.] The correction restricted to the top-$r$ principal
  directions of the between-subject, within-stage covariance estimated on the
  training subjects, for $r\in\{1,2,4\}$.
  \item[SCSN-hard.] The same estimator with one-hot pseudo-labels.
  \item[SCSN-orth.] The correction projected off the row space of the
  classifier, so that the logits are unchanged exactly.
  \item[Mismatched.] The shifts of another recording of the same fold, applied
  cyclically; the control for whether the effect requires the recording's own
  statistics.
  \item[Tent.] Entropy minimization over the affine parameters of the
  normalization layers \citep{wang2021tent}, with the published optimizer
  (SGD, momentum $0.9$, learning rate $2.5\times10^{-4}$, one pass over the
  recording), and two Adam variants (one and five steps, learning rate
  $10^{-3}$). Calibration recordings are adapted in the same way as test
  recordings.
\end{description}

The group-normalized checkpoints are evaluated with the arms that do not
require batch statistics: the uncorrected model, the post-hoc marginal
correction and SCSN. Standard AdaBN has no batch statistic to re-estimate on
these checkpoints. Standard Tent updates batch-normalization affine
parameters; a separate GroupNorm-specific entropy objective is not evaluated
here because it would introduce a new optimizer, learning-rate choice, and
step-count protocol rather than measuring the effect of stage conditioning.

\subsection{Training and splits}

Evaluation uses five subject-wise folds. Within a fold the 20\% of subjects
with the lowest hash are held out for testing, a further disjoint 20\% form
the conformal calibration set, 15\% of the remaining subjects are used for
early stopping, and the rest for training. No subject appears in more than one
of the four groups of a fold. On Sleep-EDF each fold is run with three seeds
for EEGNet and one seed for SleepCNN, giving fifteen and five runs per arm;
on HMC each fold is run with two seeds for both backbones, giving ten runs per
arm. Optimization uses Adam with learning rate
$10^{-3}$, weight decay $10^{-4}$ and cross-entropy without class weights, for
at most 40 epochs with early stopping on the validation split (patience 8).
The best validation checkpoint is kept.

\subsection{Conformal evaluation}

For each run, the calibration subjects' epochs fit a split-conformal predictor
with the cumulative (APS) non-conformity score \citep{romano2020} at
$\alpha=0.1$, and the test subjects' epochs are predicted as sets. We report
coverage, mean set size, and review workload (the fraction of epochs whose set
holds more than one stage). Because the guarantee is marginal over patients,
coverage is also computed per test subject, from which we report the
worst-covered subject, the standard deviation of per-subject coverage, the
10th percentile and the median. The calibration split is disjoint from the test
split, and the quantile is refitted after every correction.

\subsection{Identity measurements}

For every run the penultimate representation of every test epoch is extracted
under each arm, and two probes are fitted. The \emph{overall} probe is a
multinomial logistic regression from the representation to the patient index,
fitted with three-fold cross-validation over epochs. The \emph{within-stage}
probe repeats the fit separately for the epochs of each stage and averages the
five accuracies over the stages that contain at least fifty epochs. The overall
probe is sensitive to any recording-specific offset, including the one that a
marginal correction adds; the within-stage probe holds the stage fixed, so it
measures the part of the identity that is not a difference in sleep
architecture and does not react to a correction that removes only a
recording-wide mean. To keep the cost bounded on the large cohort, a fixed
random subsample of at most 5{,}000 epochs per run (3{,}000 for the
256-dimensional SleepCNN features) is used; the subsample is drawn with a fixed
seed, so every arm of a run is probed on exactly the same epochs. Chance is
$1/|\cS|$ for $|\cS|$ held-out subjects: four on Sleep-EDF and thirty-one on
HMC. Both probes use the test patients' identifiers, and the within-stage
probe additionally uses the true stage of each test epoch to form the groups.
They are diagnostics of the representation, not components of any arm: no
correction reads a test label, and the probe is fitted after the arm has
produced its representation.

\subsection{Statistics}

All comparisons between arms are paired: the arms of a run share the
checkpoint, the fold, the seed, the test epochs and the probe subsample.
Reported differences are the mean of the per-run differences, a 95\% confidence
interval from the $t$ distribution over runs, the standardised mean difference
$d_z$, and a two-sided Wilcoxon signed-rank test. Because the seeds of a fold
share their test subjects, runs within a fold are not independent; every paired
effect is therefore also reported as a sign count over folds (the sign of the
fold-averaged difference). $p$-values are unadjusted, and the number of paired
runs behind every contrast is reported so that the evidence can be weighed.
